\documentclass[10pt]{article}
\usepackage[margin=20mm]{geometry}
\usepackage{amsmath,amssymb,mathtools,fontspec,unicode-math,xurl,hyperref}
\IfFontExistsTF{DejaVu Serif}{\setmainfont{DejaVu Serif}}{\setmainfont{Latin Modern Roman}}
\setmathfont{latinmodern-math.otf}
\setlength{\parindent}{0pt}\setlength{\parskip}{6pt}\sloppy
\title{Selected rigidity arguments: initial verification}
\author{}
\date{25 October 2026}
\begin{document}
\maketitle
Abstract: This report checks selected hypotheses, reductions, and estimates in Zhenqi Jenny Wang’s arXiv:1908.10496v3. A boundary case in Corollary 1.7 is not covered by the hypotheses of the splitting theorem invoked in its proof. Its status is Needs Evidence, not a demonstrated counterexample to the rigidity conclusion. Two local algebraic/analytic checks pass. The principal splitting and iteration proofs have not been fully verified.\par
\section*{Identity, version, and scope}
The assigned paper is M25-0031, Local rigidity of weak or no hyperbolicity algebraic actions, DOI 10.1090/jams/1058. The inspected original is the 98-page author manuscript arXiv:1908.10496v3, posted 19 March 2025. All page references below are one-based PDF pages, agreeing with printed pages on the inspected pages. The PDF SHA-256 is\par
`a0aa39b077cac09ead24f9ccc58176d820fef346820361865b42cb3e2ad58ebf`.\par
No correspondence with the journal final text has been established.\par
The review traces Corollary 1.7 through Section 11.4, Theorem 9.2, and Section 3.1(11); checks the algebraic hypothesis used in Section 11.1; and recomputes the commuting-perturbation identity in Lemma 5.1 and the affine regularity recurrence in Appendix A. Definitions and statements on pp. 3, 6–7, 21–22, 27–29, 63–64, and 86–87 were read, with surrounding introductory text. This is a claim-focused, partial initial review. It does not certify Sections 6–10, Appendix C, or the paper as a whole.\par
The ambient assumptions include a compact quotient by a torsion-free cocompact irreducible lattice, the stated semisimple group hypotheses, and smooth continuous actions. Strong local rigidity permits a group automorphism of the acting group (p. 3). Geometric stability is the quantitative condition (5.4), not merely an informal stability assertion (pp. 28–29).\par
\section*{Problem 1: the scope of Corollary 1.7}
\subsection*{Source claim and dependency}
Corollary 1.7 (p. 7) is phrased for an \(\mathbb R^m\) action of upper triangular unipotents in \[G=SL(2,\mathbb R)\times SL(2,k_1)\times\cdots\times SL(2,k_n),
\qquad m\leq 1+n .\] Section 11.4 (p. 86; released `main.tex`, lines 4997–4998) obtains a splitting directly from Theorem 9.2. However, the latter uses the generating set \(E\) fixed in Section 3.1(11), p. 22: in the product case it contains \(u=E_{12}\) in the first factor and a nonzero nilpotent \(v\) from another factor. Section 9.2 (p. 64) expressly retains this choice. The definition of a friendly pair also excludes zero vectors (p. 22).\par
\subsection*{Complete boundary check}
Take \(n=1\), \(k_1=\mathbb R\), \(m=1\), and fix any quotient satisfying the ambient lattice assumptions. Define \[A=\{(I+tE_{12},I):t\in\mathbb R\}.\] Because \(E_{12}^2=0\), multiplication adds parameters, inversion negates them, and the parameter is recovered from the \((1,2)\) entry. Thus \(A\) is a closed abelian subgroup isomorphic to \(\mathbb R\), consisting of upper triangular unipotents. It obeys \(m=1\leq2\). Its Lie algebra is \[\mathop{\mathrm{Lie}}(A)=\mathbb R(E_{12},0),\qquad
\mathop{\mathrm{Lie}}(A)\cap(0\oplus\mathfrak{sl}_2(\mathbb R))=\{0\}.\] Consequently no generating set for this action contains the required nonzero second-factor vector. A change of the one-dimensional acting-group coordinate does not change that intersection. The direct application of Theorem 9.2 therefore does not cover this boundary case.\par
This calculation is a counterexample to the implication from the displayed corollary hypotheses to the hypotheses of the cited splitting theorem. It is not a counterexample to transversal local rigidity itself. In particular, no nonconjugate deformation or failure of every possible transversal family has been proved here.\par
\subsection*{Counterevidence, impact, and disposition}
The sentence immediately before Corollary 1.7 describes removing the geometric-stability assumption from Theorem 1.1. This strongly suggests that the nonzero cross-factor condition was intended to remain in force. Such an interpretation excludes the example and is material counterevidence against declaring the intended result false. Conversely, the displayed corollary does not explicitly retain that condition, and the one-sentence reduction does not supply an alternative splitting for all its written cases.\par
The judgment is Needs Evidence: resolve the intended action class and demonstrate the splitting hypotheses for that class. The affected dependency is the stated range of Corollary 1.7. No defect in Theorems 1.1 or 1.2 has been established. No complete repair, author adoption, or human acceptance is claimed.\par
\section*{Checks with no issue found within their scope}
\subsection*{Commuting perturbations}
For commuting generators \(E_i\) and perturbations \(E_i+p_i\) that still commute, \[[E_i,p_j]-[E_j,p_i]=-[p_i,p_j].\] This follows by expanding \([E_i+p_i,E_j+p_j]=0\). For each nonnegative integer \(t\), the usual product estimate obtained by Leibniz expansion and interpolation gives \[\|[p_i,p_j]\|_{C^t}\leq C_t\|p\|_{C^0}\|p\|_{C^{t+1}}.\] The interpolation step matters: the coarser displayed bound (5.2) alone is not the sharp tame estimate. In the invariant frame, integration of derivative terms against Haar measure vanishes. Averaging the linearized expression thus gives the second bound in Lemma 5.1 at \(t=0\). This verifies the local commuting-error calculation on p. 28; it does not verify the later nonlinear iteration.\par
\subsection*{The Jordan-block regularity recurrence}
Condition on the scalar estimate in Theorem A.2 (p. 87). Put \(F(t)=s_2t+s_1\). Solving the triangular system from the last coordinate upward composes \(F\) once per coordinate, so after \(j\) coordinates, \[F^{\circ j}(t)=s_2^jt+s_1\sum_{r=0}^{j-1}s_2^r .\] The base case is immediate; composing this expression with \(F\) proves the next case. This matches the recurrence (A.2) and the displayed component estimate. The independent SymPy calculation checked \(j=1,\ldots,8\) with exact zero residuals; the induction, rather than that finite calculation, establishes the formula for arbitrary finite \(j\). Theorem A.2 and its external analytic dependencies were not independently reproved. The additional last-line assertion bounding \(\lambda\) by \(\dim\mathfrak g\) was not established by this check.\par
\subsection*{The explicit rectangular subgroup}
In \(SL(6,\mathbb R)\), the nine matrices \(E_{ij}\) with \(i\in\{1,3,5\}\) and \(j\in\{2,4,6\}\) commute: both Kronecker symbols in \[[E_{ij},E_{kl}]=\delta_{jk}E_{il}-\delta_{li}E_{kj}\] vanish. They contain \(E_{12},E_{34},E_{56}\), whose three disjoint root pairs give commuting \(\mathfrak{sl}_2\) triples. Their pairwise root differences are not roots. This checks the root-space part of the reduction in Section 11.1, conditional on the geometric-stability result.\par
A scope clarification is retained separately: maximal abelian in Remark 1.5 and Corollaries 1.3–1.4 must be read consistently with the specific large-dimensional rectangular subgroup used in (3.1). Maximality by inclusion alone does not identify that subgroup. The proof’s explicit normal form is counterevidence against treating this terminology alone as a new substantive defect.\par
\section*{Conclusion and remaining work}
Item | Judgment | Repair/adoption\par
Corollary 1.7 range | Needs Evidence | No verified repair; adoption not requested\par
Lemma 5.1 local identity | No Issue Found | No repair proposed\par
Appendix A recurrence | No Issue Found, conditional | Scalar estimate not reverified\par
Rectangular-root reduction | Scope clarification | Geometric stability not reverified\par
The principal outstanding scientific work is the full verification of directional smoothing, global splitting, KAM convergence, the geometric-stability proof, and external cited results. The arXiv record inspected on 25 October 2026 lists v3 as its latest revision. A bounded title/identifier search and the author’s publication page did not identify a separate erratum; this is not a proof that none exists. No later-version or journal-final comparison was performed. No experiments are needed for the displayed algebraic checks.\par
\section*{Supporting checks}
Original pages 1, 3, 6, 7, 22, 27, 28, 29, 64, 86, and 87 were rendered locally with Poppler and personally inspected in this conversation. The first-page renderer emitted a Symbol-font warning; the version, title, author, and visible mathematical text used for identity remained legible. The other inspected renderings had no recorded stderr. Text extraction covered 98 pages, but extraction is not proof review.\par
The released TeX was inspected statically at the indicated reduction and recurrence; it was not executed. The independent script `independent\_check.py`, its exact results, extraction log, render log, and page inspection ledger are saved under the paper’s evidence directory. The symbolic package was SymPy 1.13.3. No author code was run.\par
The original compiler-absence probe preceded the authorized portable Tectonic follow-up. Its build and page inspections remain preserved as historical evidence. The current report is compiled from this exact TeX with Tectonic 0.17.0 in untrusted, cached-only mode. Exact input, dependency, log, PDF and page-image hashes, together with actual inspection of changed or reflowed pages, are recorded in the accompanying provenance. Original-review actions retain their original attribution; focused AI acceptance is recorded separately. Scientific adjudication remains pending, with no human review or author adoption claimed.\par
Backend equivalence is unverified. Human scientific adjudication remains pending.\par
9 Z. J. Wang, Local rigidity of weak or no hyperbolicity algebraic actions, arXiv:1908.10496v3 (19 March 2025), exact locally bound author manuscript. https://arxiv.org/abs/1908.10496v3. arXiv version record for 1908.10496, inspected 25 October 2026. https://arxiv.org/abs/1908.10496. Z. Wang, official Michigan State University publication page, inspected 25 October 2026. https://apps.math.msu.edu/PageSpace/pb/zwang.\par
\end{document}
